%% ============================================================================
%%  The Math and Physics of MechanicsDSL
%%  Compile with:  pdflatex MATH_AND_PHYSICS.tex   (run twice for the ToC)
%% ============================================================================

\documentclass[11pt,a4paper]{article}

\usepackage[T1]{fontenc}
\usepackage[utf8]{inputenc}
\usepackage{amsmath,amssymb}
\let\Bbbk\relax                      % amssymb and newtxmath both define \Bbbk
\usepackage{newtxtext,newtxmath}
\usepackage[margin=1.05in]{geometry}
\usepackage{booktabs}
\usepackage{tabularx}
\usepackage{array}
\usepackage[expansion=false]{microtype}
\usepackage{enumitem}
\usepackage{listings}
\usepackage{xcolor}
\usepackage{titlesec}
\usepackage[colorlinks=true,linkcolor=linknavy,urlcolor=linknavy,citecolor=linknavy]{hyperref}

\definecolor{linknavy}{RGB}{20,60,110}
\definecolor{codebg}{RGB}{247,247,244}
\definecolor{coderule}{RGB}{205,205,198}
\definecolor{codekey}{RGB}{25,80,140}

%% --- Sectioning -------------------------------------------------------------
\titleformat{\section}{\normalfont\Large\bfseries}{\thesection}{0.8em}{}
\titleformat{\subsection}{\normalfont\large\bfseries}{\thesubsection}{0.7em}{}
\titlespacing*{\section}{0pt}{2.2ex plus 1ex minus .2ex}{1.2ex plus .2ex}

%% --- Tables -----------------------------------------------------------------
\renewcommand{\arraystretch}{1.25}
\newcolumntype{L}[1]{>{\raggedright\arraybackslash}p{#1}}

%% --- Code listings ----------------------------------------------------------
\lstdefinestyle{dsl}{
  basicstyle=\ttfamily\small,
  keywordstyle=\color{codekey}\bfseries,
  backgroundcolor=\color{codebg},
  frame=single,
  rulecolor=\color{coderule},
  framesep=6pt,
  xleftmargin=10pt,
  xrightmargin=10pt,
  breaklines=true,
  columns=fullflexible,
  showstringspaces=false,
  morekeywords={system,defvar,parameter,lagrangian,initial,force,rayleigh,constraint,hamiltonian}
}
\lstset{style=dsl}

%% --- Convenience macros -----------------------------------------------------
\newcommand{\ud}{\mathrm{d}}
\newcommand{\dd}[2]{\frac{\partial #1}{\partial #2}}
\newcommand{\ddt}[1]{\frac{\ud #1}{\ud t}}
\newcommand{\code}[1]{\texttt{#1}}

\title{\textbf{The Math and Physics of MechanicsDSL}}
\author{}
\date{}

% ============================================================================
\begin{document}
% ============================================================================

\maketitle
\vspace{-3.2em}
\begin{center}
\large\itshape A complete account of what this library computes,\\ and why it is built the way it is.
\end{center}
\vspace{1.4em}

\hrule
\vspace{1.2em}

\section*{How to read this document}
\addcontentsline{toc}{section}{How to read this document}

No prior calculus is assumed. Section~\ref{sec:notation} establishes the notation, and unfamiliar
constructions are explained where they first appear. Readers already familiar with Lagrangian
mechanics may proceed directly to Section~\ref{sec:engine}.

The document divides in two:

\begin{itemize}[leftmargin=1.4em,itemsep=0.3em]
  \item \textbf{Sections \ref{sec:notation}--\ref{sec:energy}} describe the core engine and are
        intended to be read in order.
  \item \textbf{Section \ref{sec:domains}} is a reference catalog of the sixteen physics domains
        the library implements. Each entry opens with a one-sentence statement of scope; the
        material beneath documents coverage rather than inviting continuous reading.
\end{itemize}

\vspace{0.8em}
\hrule
\vspace{0.6em}

\tableofcontents

\vspace{1.2em}
\hrule

% ============================================================================
\section{Notation}
\label{sec:notation}
% ============================================================================

\textbf{Symbols denote quantities.} Here $m$ is mass, $l$ length, $g$ gravitational acceleration
($9.81\ \mathrm{m/s^2}$ at Earth's surface), $t$ time, and $E$ energy. Greek letters follow the same
convention: $\theta$ typically an angle, $\omega$ an angular frequency.

\medskip
\noindent\textbf{$q$ denotes a generalized coordinate} --- a single number specifying the
configuration of a system. For a pendulum this might be the angle from vertical; for a block on a
track, the distance along it. Where several are required they are indexed $q_1, q_2, \dots$. The
term \emph{generalized} signals that no particular kind of coordinate is assumed: an angle and a
distance are treated identically.

\medskip
\noindent\textbf{An overdot denotes differentiation with respect to time} --- the rate at which a
quantity changes, per second.

\begin{center}
\begin{tabular}{@{}clL{7.6cm}@{}}
\toprule
\textbf{Symbol} & \textbf{Read as} & \textbf{Quantity} \\
\midrule
$q$        & ``$q$''              & Position \\
$\dot{q}$  & ``$q$-dot''          & Velocity --- the rate of change of position \\
$\ddot{q}$ & ``$q$-double-dot''   & Acceleration --- the rate of change of velocity \\
\bottomrule
\end{tabular}
\end{center}

\noindent For a pendulum with angular coordinate $\theta$, then, $\dot\theta$ is its angular
velocity and $\ddot\theta$ its angular acceleration. This is the entire convention.

\medskip
\noindent\textbf{The partial derivative $\partial f/\partial x$} measures how a quantity $f$
responds to variation in $x$ while every other variable is held fixed. On a topographic map,
$\partial(\text{elevation})/\partial(\text{easting})$ is the gradient one would measure walking due
east; variation in the north--south direction contributes nothing. The notation isolates a single
direction of dependence.

\medskip
\noindent\textbf{The operator $\ud/\ud t$} denotes the total rate of change with respect to time,
accounting for every quantity that varies as time passes.

\medskip
\noindent\textbf{Additional notation:}

\begin{center}
\begin{tabular}{@{}cL{9.6cm}@{}}
\toprule
\textbf{Symbol} & \textbf{Meaning} \\
\midrule
$\sum$        & Sum over the indicated terms \\
$\int$        & Integral --- a continuous sum \\
$\nabla$      & Gradient --- direction and magnitude of steepest increase \\
$\approx$     & Approximately equal \\
$\delta_{ij}$ & Equal to $1$ when $i = j$, otherwise $0$ \\
\bottomrule
\end{tabular}
\end{center}

\noindent Nothing beyond this is assumed.

% ============================================================================
\section{What the library does}
\label{sec:what}
% ============================================================================

A physical system is specified by a single expression --- its \textbf{Lagrangian}. From that
expression the library derives the equations of motion, integrates them numerically, and can
translate them into thirteen target languages.

A complete pendulum specification:

\begin{lstlisting}
\system{pendulum}
\defvar{theta}{Angle}{rad}
\parameter{m}{1.0}{kg}
\parameter{l}{1.0}{m}
\parameter{g}{9.81}{m/s^2}
\lagrangian{\frac{1}{2}*m*l^2*\dot{theta}^2 - m*g*l*(1 - \cos{theta})}
\initial{theta=0.5, theta_dot=0.0}
\end{lstlisting}

\noindent The declarations name the system, identify \code{theta} as the generalized coordinate,
fix the physical parameters, state the Lagrangian, and set initial conditions. Everything
downstream --- equations of motion, energy expressions, generated C\texttt{++} --- follows from the
Lagrangian line alone. No per-system code is written by hand.

The pipeline comprises five stages:

\begin{center}
\begin{tabular}{@{}lL{7.4cm}l@{}}
\toprule
\textbf{Stage} & \textbf{Function} & \textbf{Module} \\
\midrule
Parse      & Source text to abstract syntax tree        & \code{parser/core.py} \\
Symbolize  & Syntax tree to symbolic algebra            & \code{symbolic.py} \\
Derive     & Lagrangian to equations of motion          & \code{symbolic.py} \\
Solve      & Equations to numerical trajectory          & \code{solver/} \\
Emit       & Equations to source in a target language   & \code{codegen/} \\
\bottomrule
\end{tabular}
\end{center}

% ============================================================================
\section{Energy formulations versus force formulations}
\label{sec:formulations}
% ============================================================================

\subsection{The two approaches}

Classical mechanics admits two equivalent formulations. They yield identical predictions and differ
in what they require of the person applying them.

\medskip
\noindent\textbf{The Newtonian formulation} works with forces:
\begin{equation}
\vec{F} = m\vec{a}
\end{equation}
Applying it requires identifying every force acting on every body and resolving each into components
along chosen axes. The bookkeeping is vectorial, and it is the user's responsibility: an omitted
force yields a wrong answer with no indication that anything is missing.

\medskip
\noindent\textbf{The Lagrangian formulation} works with energy. Define a single scalar quantity, the
Lagrangian:
\begin{equation}
L = T - V
\end{equation}
where $T$ is the kinetic energy --- the energy associated with motion --- and $V$ is the potential
energy, the energy associated with configuration. A raised mass and a compressed spring both store
potential energy; both convert it to kinetic energy when released.

$L$ is a scalar. It has no direction and no components.

\subsection{Why the difference rather than the sum}
\label{sec:action}

The subtraction is not a convention chosen for convenience; it is fixed by the variational principle
underlying the formulation.

Among all conceivable paths a system might follow between two configurations, the physical path is
distinguished by a property of the \textbf{action} --- the integral of $L$ along the path. The
physical path is the one for which the action is \textbf{stationary}: perturb the path
infinitesimally in any direction and the action is unchanged to first order.

The analogy is to a minimum of a function. A ball at the base of a bowl sits where displacement in
any direction produces no first-order change in height. The physical trajectory occupies the
analogous position within the space of all possible trajectories.

Taking $T - V$ is what makes the stationary-action condition reproduce Newtonian dynamics. The
result is the Euler--Lagrange equation, derived next.

Using the library requires no commitment to this justification. It does explain the minus sign.

\subsection{The Euler--Lagrange equation}

For each generalized coordinate $q$, the equation of motion is
\begin{equation}
\ddt{}\left(\dd{L}{\dot q}\right) - \dd{L}{q} = 0
\label{eq:el}
\end{equation}

\noindent Read from the inside outward, this prescribes three operations:

\begin{enumerate}[leftmargin=1.8em,itemsep=0.25em]
  \item $\partial L/\partial \dot q$ --- differentiate $L$ with respect to the velocity.
  \item $\ud/\ud t\,(\cdot)$ --- differentiate that result with respect to time.
  \item $\partial L/\partial q$ --- differentiate $L$ with respect to the coordinate itself, and
        subtract.
\end{enumerate}

\noindent The procedure is purely mechanical and identical for every system and every coordinate.
That invariance is what makes it automatable, and automating it is what this library is.

\subsection{Worked example: the simple pendulum}
\label{sec:worked}

Applying the procedure once, in full, to the system specified in Section~\ref{sec:what}.

\medskip
\noindent\textbf{Setup.} A point mass $m$ on a rigid massless rod of length $l$, displaced by angle
$\theta$ from vertical.

The mass travels on a circular arc of radius $l$, so its speed is $l\dot\theta$. Kinetic energy is
$\tfrac{1}{2}mv^2$:
\begin{equation}
T = \tfrac{1}{2}m l^2\dot\theta^2
\end{equation}

Displacement through angle $\theta$ raises the mass by $l(1 - \cos\theta)$, giving gravitational
potential energy $mgh$:
\begin{equation}
V = mgl\,(1 - \cos\theta)
\end{equation}

The Lagrangian is their difference:
\begin{equation}
L = \tfrac{1}{2}m l^2\dot\theta^2 - mgl(1 - \cos\theta)
\end{equation}

\medskip
\noindent\textbf{Step 1: differentiate with respect to $\dot\theta$.}
Only the kinetic term depends on $\dot\theta$. Differentiating a squared quantity $x^2$ yields $2x$,
so $\tfrac{1}{2}ml^2\dot\theta^2$ yields $\tfrac{1}{2}ml^2 \cdot 2\dot\theta$:
\begin{equation}
\dd{L}{\dot\theta} = m l^2 \dot\theta
\end{equation}
This quantity is the pendulum's \textbf{angular momentum}. It emerges from the procedure without
having been sought.

\medskip
\noindent\textbf{Step 2: differentiate with respect to time.}
Here $m$ and $l$ are constants, so only $\dot\theta$ varies; its time derivative is $\ddot\theta$:
\begin{equation}
\ddt{}\left(m l^2 \dot\theta\right) = m l^2 \ddot\theta
\end{equation}

\medskip
\noindent\textbf{Step 3: differentiate with respect to $\theta$ and subtract.}
Only the potential term depends on $\theta$. Expanding,
$-mgl(1 - \cos\theta) = -mgl + mgl\cos\theta$; the constant contributes nothing, and the derivative
of $\cos\theta$ is $-\sin\theta$:
\begin{equation}
\dd{L}{\theta} = -mgl\sin\theta
\end{equation}

Assembling:
\begin{equation}
m l^2 \ddot\theta + mgl\sin\theta = 0
\end{equation}

Dividing by $ml^2$:
\begin{equation}
\boxed{\ \ddot\theta = -\frac{g}{l}\sin\theta\ }
\end{equation}

\medskip
\noindent\textbf{Interpretation.} Angular acceleration is proportional to $-\sin\theta$. The
negative sign gives a restoring effect --- acceleration is always directed back toward vertical ---
and the magnitude grows with displacement. Note that $m$ has cancelled: the period is independent of
the mass, as observed.

\medskip
\noindent\textbf{Verification.} For small displacements $\sin\theta \approx \theta$, reducing the
equation to $\ddot\theta = -(g/l)\theta$, the equation of simple harmonic motion with angular
frequency $\omega = \sqrt{g/l}$ and period
\begin{equation}
T = 2\pi\sqrt{l/g}
\end{equation}
which is the standard result. It follows from the definition of $T - V$ and three differentiations.

\medskip
\noindent\emph{No force was identified at any point, and the rod's tension never appeared.} That
observation motivates the next section.

% ============================================================================
\section{Why the Lagrangian formulation}
\label{sec:why}
% ============================================================================

Both formulations are correct. The question is which can be carried out by software on the user's
behalf. Six considerations follow, then the cases where the choice is a liability.

\subsection{A scalar expression is representable as text}

$L = T - V$ is a single algebraic expression. It can be tokenized, parsed, differentiated
symbolically, and simplified by a computer algebra system.

The Newtonian formulation requires that the analysis already be done: forces enumerated, components
resolved, frame selected. There is no natural textual representation of ``the tension in the rod''
from which software could infer the intended physics. There is a natural textual representation of
kinetic energy.

\subsection{Constraint forces do not appear}
\label{sec:noconstraintforce}

The rod in Section~\ref{sec:worked} exerts a tension on the mass. A Newtonian treatment must account
for it --- despite its doing no work, having no influence on the period, and being irrelevant to the
question asked. It would be solved for and then eliminated.

In Section~\ref{sec:worked} it never entered the calculation. Because the configuration was
described by the single quantity that can actually vary, the constraint was incorporated from the
outset: motion the rod forbids is not expressible in the coordinate.

The advantage compounds with system size. A double pendulum in Cartesian coordinates requires four
position variables, two rigid-length constraints, and two unknown tensions. In generalized
coordinates it requires two angles.

\subsection{The formulation is coordinate-independent}

The Euler--Lagrange equation takes the same form in every coordinate system --- Cartesian, polar,
spherical, or arbitrary curvilinear.

Newton's second law does not. Transforming to polar coordinates introduces centrifugal and Coriolis
terms that must be derived separately and included by hand; omitting one produces an error that is
difficult to detect.

A domain-specific language cannot anticipate which coordinates a user will choose, so it requires a
formulation indifferent to that choice. This is why \code{\textbackslash defvar\{theta\}} and
\code{\textbackslash defvar\{x\}} are handled identically by the parser --- no stage downstream needs
to distinguish them.

\subsection{Coupling is derived rather than assembled}
\label{sec:coupling}

In a double pendulum each arm's angular acceleration depends on the other's. Extracting that
dependence from a force analysis is lengthy and error-prone.

From the Lagrangian it is automatic. A coupling term such as $\cos(\theta_1 - \theta_2)$ propagates
correctly through the differentiation steps with no special handling.

This is why \code{derive\_equations\_of\_motion} in \code{symbolic.py} promotes \emph{all}
coordinates to functions of time simultaneously before differentiating. Treating them sequentially
would drop the cross-terms and produce equations that are plausible and wrong.

\subsection{The structures required downstream are already present}

Available directly from the Lagrangian at no additional cost:

\begin{itemize}[leftmargin=1.4em,itemsep=0.25em]
  \item \textbf{Total energy}, $T + V$ --- and the code recovers $T$ and $V$ from $L$ automatically
        (Section~\ref{sec:codegen}).
  \item \textbf{Conjugate momentum}, $\partial L/\partial\dot q$ --- obtained incidentally in Step 1
        of Section~\ref{sec:worked}.
  \item \textbf{Conservation laws}, from the symmetries of $L$ (Section~\ref{sec:noether}).
  \item \textbf{The symplectic structure} required by the integrators that prevent energy drift
        (Section~\ref{sec:symplectic}).
\end{itemize}

A set of force vectors carries none of this. Each would require separate implementation per system.

\subsection{The formulation generalizes}

The same relation~\eqref{eq:el} applies to charged particles in electromagnetic fields, to
relativistic particles, to continuous fields, and to geodesic motion in curved spacetime. This is how
the \code{domains/} package spans sixteen fields on one engine. Newton's second law does not extend
in this way.

\subsection{Where the choice is a liability}

The formulation is the weaker option when:

\begin{itemize}[leftmargin=1.4em,itemsep=0.3em]
  \item \textbf{A constraint force is the quantity of interest.} If the question is whether the rod
        will fail, the tension is required --- and recovering what the formulation eliminated
        requires the Lagrange multiplier machinery of Section~\ref{sec:constraints}.
  \item \textbf{Non-conservative forces are present.} Friction and drag dissipate energy and admit no
        potential, so they cannot be expressed within $T - V$ and must be reintroduced explicitly.
        The DSL provides \code{\textbackslash force\{\}} and \code{\textbackslash rayleigh\{\}} for
        this purpose (Section~\ref{sec:dissipation}).
  \item \textbf{The system is trivial.} For a single particle in rectilinear motion, $F = ma$ is more
        direct.
\end{itemize}

For the remaining majority of problems the scalar formulation is preferable because it can be
mechanized --- and mechanization is the purpose of a compiler.

% ============================================================================
\section{The symbolic engine}
\label{sec:engine}
% ============================================================================

\subsection{Deriving the equations}

For each coordinate the engine executes the procedure of Section~\ref{sec:worked} symbolically:

\begin{enumerate}[leftmargin=1.8em,itemsep=0.2em]
  \item Promote every coordinate to a function of time --- \emph{all coordinates simultaneously},
        preserving coupling (Section~\ref{sec:coupling}).
  \item Differentiate with respect to velocity, then with respect to time.
  \item Subtract the derivative with respect to position.
  \item Substitute back to plain symbols and simplify.
\end{enumerate}

Simplification runs under a timeout, since symbolic simplification has no useful worst-case bound. On
expiry the unsimplified expression is used --- less readable, equally correct.

\subsection{Solving for accelerations: the mass matrix}
\label{sec:massmatrix}

The Euler--Lagrange equations are always \textbf{linear in the accelerations}: each $\ddot q$ appears
to the first power, never squared or inside a transcendental function. The system therefore admits
the compact form
\begin{equation}
M(q)\,\ddot{q} = F(q, \dot q, t)
\end{equation}

\begin{itemize}[leftmargin=1.4em,itemsep=0.3em]
  \item $M$ is the \textbf{mass matrix}, an $n \times n$ array whose entry $M_{ij}$ is the
        coefficient of $\ddot q_j$ in equation $i$. Diagonal entries express each coordinate's own
        inertia; off-diagonal entries express inertial coupling --- the extent to which accelerating
        one coordinate exerts an inertial effect on another. For a single particle it reduces to the
        mass.
  \item $F$ is the \textbf{generalized force vector}, comprising every term without an acceleration.
\end{itemize}

Solving requires inverting $M$. Two paths are implemented:

\begin{itemize}[leftmargin=1.4em,itemsep=0.25em]
  \item \textbf{One coordinate:} direct solution. Given $A\ddot q + B = 0$, then $\ddot q = -B/A$.
  \item \textbf{Several coordinates:} simultaneous solution of the full linear system, with explicit
        matrix inversion as a fallback. Sequential solution would be invalid, since each acceleration
        appears in the others' equations.
\end{itemize}

If $\det M = 0$ the mass matrix is \textbf{singular} and the system has no unique solution.
Physically this indicates a direction of motion carrying no inertia --- a degenerate Lagrangian. The
condition is detected and reported rather than propagated.

\subsection{The Hamiltonian formulation}
\label{sec:hamiltonian}

An alternative packaging of the same dynamics replaces velocity with momentum as the second state
variable. The \textbf{conjugate momentum} is
\begin{equation}
p = \dd{L}{\dot q}
\end{equation}
which for an unconstrained particle reduces to $mv$, and for the pendulum of
Section~\ref{sec:worked} gave $ml^2\dot\theta$.

The change of variables is a \textbf{Legendre transform} --- a systematic method for re-expressing a
function in terms of its own derivative --- and yields the \textbf{Hamiltonian}:
\begin{equation}
H = \sum_i p_i \dot q_i - L
\end{equation}
with every velocity eliminated in favour of momentum. For systems whose kinetic energy is quadratic
in velocity, $H = T + V$: the total energy.

The dynamics then take the symmetric first-order form of \textbf{Hamilton's equations}:
\begin{equation}
\dot q_i = \dd{H}{p_i}, \qquad \dot p_i = -\dd{H}{q_i}
\end{equation}

The motivation is geometric. The space of states $(q, p)$ --- \textbf{phase space}, in which a single
point specifies the complete instantaneous state --- carries a structure (the symplectic form) that
certain integrators preserve exactly. Preserving it is what bounds long-term energy error
(Section~\ref{sec:symplectic}).

\subsection{Constraints and Lagrange multipliers}
\label{sec:constraints}

A \textbf{holonomic constraint} restricts configuration alone and is written $g(q) = 0$ --- a bead
confined to a wire, or two masses at fixed separation.

It is imposed by augmenting the Lagrangian with the constraint and an undetermined coefficient:
\begin{equation}
L' = L + \sum_k \lambda_k\, g_k(q)
\end{equation}

Each multiplier $\lambda_k$ is an additional unknown. Solving for it yields more than the constraint:
$\lambda$ is proportional to the constraint force, which is how the tension eliminated in
Section~\ref{sec:noconstraintforce} is recovered when it is actually wanted.

One point requires care in implementation. The constraint must be differentiated \emph{twice} with
respect to time:
\begin{equation}
g = 0 \implies \dot g = 0 \implies \ddot g = 0
\end{equation}
If a separation is constant, its rate of change vanishes, as does the rate of change of that rate.
Only the second derivative contains accelerations --- the unknowns being solved for. Stopping at the
first leaves the multipliers undetermined, and they propagate into the accelerations as unbound
symbols. Accelerations and multipliers are therefore solved as a single combined linear system.

A \textbf{non-holonomic constraint} restricts velocities in a form that cannot be integrated to a
configuration constraint:
\begin{equation}
\sum_j a_{ij}(q)\,\dot q_j + b_i(q,t) = 0
\end{equation}
Rolling without slipping is the standard case: the constraint restricts instantaneous velocity while
leaving every position on the plane reachable. These are handled by the d'Alembert--Lagrange and
Appell formulations.

The \textbf{constraint Jacobian} $\partial g/\partial q$ is rank-checked. Rank below the number of
constraints indicates redundancy --- the same restriction expressed more than once --- which makes
the augmented system unsolvable.

\subsection{Noether's theorem}
\label{sec:noether}

\emph{Every continuous symmetry of the action corresponds to a conserved quantity.}

\begin{center}
\begin{tabular}{@{}L{7.2cm}L{5.4cm}@{}}
\toprule
\textbf{Symmetry} & \textbf{Conserved quantity} \\
\midrule
Time translation ($L$ has no explicit $t$) & Energy \\
Spatial translation                        & Linear momentum \\
Rotation                                   & Angular momentum \\
Galilean boost                             & Centre-of-mass motion \\
\bottomrule
\end{tabular}
\end{center}

Detection is direct. If a coordinate $q$ is absent from $L$ while its velocity $\dot q$ is present,
then $\partial L/\partial q = 0$ and the Euler--Lagrange equation reduces to
\begin{equation}
\ddt{}\left(\dd{L}{\dot q}\right) = 0
\end{equation}
so the conjugate momentum is constant. Such a coordinate is termed \textbf{cyclic}, and the library
identifies conservation laws by scanning for them.

For a particle in a central potential the orbital angle is cyclic --- the potential is rotationally
symmetric --- so angular momentum is conserved. This is Kepler's equal-areas law, obtained by
inspecting which symbols occur in an expression.

\subsection{Dissipation and applied forces}
\label{sec:dissipation}

Dissipative forces admit no potential and are introduced as additional terms:
\begin{equation}
\ddt{}\left(\dd{L}{\dot q_i}\right) - \dd{L}{q_i} + \dd{\mathcal{F}}{\dot q_i} = Q_i
\end{equation}

Here $\mathcal{F}$ is the \textbf{Rayleigh dissipation function}, typically $\tfrac{1}{2}b\dot q^2$
for viscous damping. The associated generalized force is $-\partial\mathcal{F}/\partial\dot q_i$,
whose sign guarantees opposition to the motion. The term $Q_i$ denotes externally applied generalized
forces --- actuators, applied torques.

Implemented friction models:

\begin{center}
\begin{tabular}{@{}lcL{6.6cm}@{}}
\toprule
\textbf{Model} & \textbf{Form} & \textbf{Regime} \\
\midrule
Viscous  & $F = -bv$                      & Resistance proportional to speed --- fluids \\
Coulomb  & $F = -\mu N \,\mathrm{sgn}(v)$ & Constant magnitude opposing motion --- dry sliding \\
Stribeck & Composite                      & Static, Coulomb and viscous combined; reproduces stick--slip \\
\bottomrule
\end{tabular}
\end{center}

\subsection{Small oscillations and normal modes}

Expanded to second order about a stable equilibrium, any system reduces to
\begin{equation}
M\ddot q + Kq = 0
\end{equation}
where $M$ is the mass matrix of Section~\ref{sec:massmatrix} and $K$ is the \textbf{stiffness
matrix}, $K_{ij} = \partial^2 V/\partial q_i \partial q_j$ --- the curvature of the potential at the
equilibrium.

The \textbf{normal modes} are the configurations in which the entire system oscillates at a single
frequency with fixed relative amplitudes. They are obtained from the \textbf{generalized eigenvalue
problem}
\begin{equation}
K\phi = \omega^2 M \phi
\end{equation}
whose eigenvalues give the natural frequencies $\omega$ and whose eigenvectors $\phi$ give the mode
shapes.

Two pendulums coupled by a spring admit two modes: in-phase and out-of-phase oscillation. Any motion
of the system is a superposition of these. The reduction of arbitrary coupled motion to a sum of
independent oscillations is the reason normal-mode analysis is central to structural and molecular
dynamics.

\subsection{Stability and chaos}

Equilibria are located by setting all derivatives to zero. The system is then linearized about each,
and the eigenvalues of the resulting Jacobian classify it:

\begin{center}
\begin{tabular}{@{}L{6.2cm}L{6.4cm}@{}}
\toprule
\textbf{Eigenvalues} & \textbf{Classification} \\
\midrule
All real parts negative & \textbf{Stable} --- perturbations decay \\
All purely imaginary    & \textbf{Marginally stable} --- perturbations persist without growth \\
Any positive real part  & \textbf{Unstable} --- perturbations grow \\
\bottomrule
\end{tabular}
\end{center}

For chaotic systems a sharper measure applies. The \textbf{largest Lyapunov exponent} $\lambda$
characterizes the divergence of initially neighbouring trajectories:
\begin{equation}
\delta(t) \approx \delta_0\, e^{\lambda t}
\end{equation}
A positive $\lambda$ implies exponential sensitivity to initial conditions and a prediction horizon
that degrades logarithmically with measurement precision. This is the operational definition of
chaos, and it is why a double pendulum released twice from nominally identical states diverges within
seconds. The exponent is estimated from trajectory data by Wolf's algorithm with periodic
renormalization.

\subsection{Singularity detection}

Derived equations are analyzed before integration for conditions that will fail numerically:

\begin{itemize}[leftmargin=1.4em,itemsep=0.25em]
  \item \textbf{Division by zero} --- denominators possessing real roots.
  \item \textbf{Trigonometric singularities} --- $\tan$ at $\pi/2$, $\cot$ at $0$.
  \item \textbf{Singular mass matrix} --- configurations at which $\det M = 0$
        (Section~\ref{sec:massmatrix}).
  \item \textbf{Gimbal lock} --- in three-angle orientation representations, configurations where two
        rotation axes become parallel and a degree of freedom is lost. The signature is a
        $1/\cos\theta$ factor. The remedy is a quaternion representation
        (Section~\ref{sec:rigidbody}).
  \item \textbf{Dependent constraints} --- a rank-deficient constraint Jacobian.
  \item \textbf{Unstable equilibria} --- stationary points at a maximum rather than a minimum of the
        potential.
\end{itemize}

Each warning reports the condition, the coordinates involved, and a suggested remedy.

\subsection{Dimensional analysis}

Units are represented as exponent vectors over base dimensions together with a scale factor.
Multiplication adds exponents; division subtracts them.

This detects dimensionally inconsistent expressions --- a length added to a time --- at specification
time rather than allowing them to reach the solver, where they would produce a numerically valid and
physically meaningless result.

% ============================================================================
\section{Numerical integration}
\label{sec:numerics}
% ============================================================================

An equation of motion gives the instantaneous acceleration. Obtaining a trajectory requires advancing
the state in discrete steps, typically many thousands of them. The choice of stepping scheme
materially affects the result.

\subsection{General-purpose integration}

The state vector interleaves positions and velocities as
$[q_1, \dot q_1, q_2, \dot q_2, \dots]$, converting $n$ second-order equations into $2n$ first-order
equations suitable for standard solvers.

Available methods: \textbf{RK45} (adaptive Runge--Kutta, the default), \textbf{RK23},
\textbf{DOP853} (eighth order), \textbf{Radau} and \textbf{BDF} (implicit), and \textbf{LSODA}
(automatic switching).

\textbf{Stiffness} warrants explanation. A system is stiff when it contains dynamics on widely
separated timescales --- a fast oscillation superimposed on a slow evolution. Explicit methods must
resolve the fastest timescale for stability reasons, even during intervals when it is dynamically
irrelevant, making them prohibitively slow. Implicit methods remain stable at larger steps. The
library performs a brief trial integration to detect stiffness and selects a method based on system
size and integration span.

\subsection{Symplectic integrators}
\label{sec:symplectic}

Conventional methods exhibit \textbf{energy drift}. Per-step truncation errors accumulate with a
consistent bias, and over long integrations a bound orbit spirals inward or outward for reasons that
are entirely numerical.

\textbf{Symplectic integrators} eliminate this. Rather than minimizing local truncation error, they
exactly preserve the geometric structure of phase space (Section~\ref{sec:hamiltonian}). The
consequence is that energy error oscillates within a bounded envelope indefinitely rather than
accumulating --- the property required for long-duration celestial mechanics.

They operate on separable Hamiltonians $H(q,p) = T(p) + V(q)$ by composing two exactly integrable
sub-flows:

\begin{itemize}[leftmargin=1.4em,itemsep=0.2em]
  \item a \textbf{kick}, advancing momentum under the potential, $p \to p - h\nabla V(q)$;
  \item a \textbf{drift}, advancing position under the kinetic term, $q \to q + h\nabla T(p)$.
\end{itemize}

Higher order is achieved by composing these with coefficients chosen so that low-order error terms
cancel.

\begin{center}
\begin{tabular}{@{}lcL{7.0cm}@{}}
\toprule
\textbf{Method} & \textbf{Order} & \textbf{Characteristics} \\
\midrule
St\"ormer--Verlet / Leapfrog & 2    & Explicit, time-reversible; the standard choice \\
Ruth-3                       & 3    & Minimal stage count \\
Yoshida-4                    & 4    & Triple-jump composition,
                                      $c_1 = 1/(2 - 2^{1/3})$, $c_2 = -2^{1/3}/(2 - 2^{1/3})$ \\
McLachlan-4                  & 4    & Five stages, reduced error constant \\
PEFRL                        & 4    & Optimized coefficients; best long-term energy behaviour in
                                      this class \\
Yoshida-6, Yoshida-8         & 6, 8 & High-precision celestial mechanics \\
\bottomrule
\end{tabular}
\end{center}

Several compositions carry \textbf{negative coefficients}, meaning the integrator advances backward in
time during intermediate stages. This is deliberate: the retrograde stage is what cancels the leading
error terms.

An \textbf{energy drift monitor} records initial and instantaneous energy throughout a run, reporting
maximum relative error and peak-to-peak oscillation, so that conservation can be verified rather than
assumed.

\subsection{Variational integrators}

A structurally different approach: rather than discretizing the equations of motion, discretize the
variational principle of Section~\ref{sec:action} directly. The action integral becomes a sum over
path segments,
\begin{equation}
S_d = \sum_k L_d(q_k, q_{k+1}, h)
\end{equation}
and requiring it to be stationary yields the \textbf{discrete Euler--Lagrange equation}
\begin{equation}
D_2 L_d(q_{k-1}, q_k) + D_1 L_d(q_k, q_{k+1}) = 0
\end{equation}
solved for $q_{k+1}$ at each step by Newton iteration.

Because the scheme derives from a genuine variational principle, it inherits a discrete Noether
theorem: momentum maps are conserved \emph{exactly}, to machine precision, rather than approximately.

\begin{center}
\begin{tabular}{@{}L{7.6cm}c@{}}
\toprule
\textbf{Discretization of $L_d$} & \textbf{Order} \\
\midrule
Midpoint rule                            & 2 \\
Trapezoidal rule                         & 2 \\
Galerkin, $s$ Gauss--Legendre points     & $2s$ \\
\bottomrule
\end{tabular}
\end{center}

The per-step nonlinear solve makes these more expensive than explicit methods, in exchange for exact
structural conservation.

\subsection{Generated code}
\label{sec:codegen}

The C\texttt{++}, Rust, Julia, Fortran, CUDA, MATLAB, JavaScript, WebAssembly, Arduino and ARM
backends emit a fixed-step \textbf{fourth-order Runge--Kutta} integrator:
\begin{equation}
y_{n+1} = y_n + \tfrac{1}{6}\left(k_1 + 2k_2 + 2k_3 + k_4\right)
\end{equation}
with $k_1$ evaluated at the interval start, $k_2$ and $k_3$ at the midpoint, and $k_4$ at the end. It
is fourth-order accurate and entirely self-contained --- the relevant property for a target with no
runtime library. The Python backend calls SciPy instead, since it is available.

Generated code also emits an energy function. Separating $T$ and $V$ from $L$ uses \textbf{Euler's
theorem on homogeneous functions}: kinetic energy is quadratic in the velocities, so it satisfies
$\sum \dot q\,\partial L/\partial \dot q = 2T$, which identifies it uniquely; $V$ is the remainder.

\subsection{Inverse problems}

Forward simulation maps parameters to trajectories. The inverse problem maps observed trajectories
back to parameters --- recovering a physical pendulum's length from recorded motion, for instance.

\begin{itemize}[leftmargin=1.4em,itemsep=0.3em]
  \item \textbf{Parameter estimation} minimizes the residual between simulation and observation,
        using local optimization (fast, requires a reasonable initial guess) or differential
        evolution (global, no guess required). Returns covariance and confidence intervals rather
        than point estimates alone.
  \item \textbf{Sensitivity analysis} computes \textbf{Sobol indices}, which decompose output
        variance among the inputs. First-order indices give each parameter's independent
        contribution; total-order indices additionally capture its interactions.
  \item \textbf{Uncertainty quantification} uses MCMC to sample the posterior distribution over
        parameters, yielding credible intervals, and bootstrap resampling as an independent check.
\end{itemize}

% ============================================================================
\section{Energy as a correctness criterion}
\label{sec:energy}
% ============================================================================

Energy conservation is the primary validation check available, because a simulation that dissipates
energy numerically produces trajectories that appear plausible and are wrong.

For standard systems the library uses closed-form expressions.

\medskip
\noindent\textbf{Simple pendulum.}
\begin{equation}
T = \tfrac{1}{2}ml^2\dot\theta^2, \qquad V = -mgl\cos\theta
\end{equation}

\noindent\textbf{Double pendulum.}
\begin{align}
T &= \tfrac{1}{2}m_1 l_1^2\dot\theta_1^2
   + \tfrac{1}{2}m_2\!\left(l_1^2\dot\theta_1^2 + l_2^2\dot\theta_2^2
   + 2l_1 l_2 \dot\theta_1\dot\theta_2\cos(\theta_1 - \theta_2)\right) \\
V &= -m_1 g l_1\cos\theta_1 - m_2 g\left(l_1\cos\theta_1 + l_2\cos\theta_2\right)
\end{align}

\noindent\textbf{Spherical pendulum.}
\begin{equation}
T = \tfrac{1}{2}ml^2\!\left(\dot\theta^2 + \sin^2\!\theta\,\dot\phi^2\right), \qquad
V = mgl(1 - \cos\theta)
\end{equation}

\noindent\textbf{Harmonic oscillator.}
\begin{equation}
T = \tfrac{1}{2}mv^2, \qquad V = \tfrac{1}{2}kx^2
\end{equation}

\medskip
The double pendulum's cross term $2l_1 l_2\dot\theta_1\dot\theta_2\cos(\theta_1 - \theta_2)$ is the
inertial coupling between the arms, and it is the source of the system's chaotic behaviour. It
follows immediately from writing down the kinetic energy; deriving it from a force analysis requires
substantially more work.

Potential energies are offset so that the minimum reads zero --- a pendulum at rest has $V = 0$. This
affects only the plotted baseline, since potential energy is defined up to an additive constant and
only differences are physically meaningful.

% ============================================================================
\section{Domain reference}
\label{sec:domains}
% ============================================================================

The \code{domains/} package implements sixteen fields beyond the core engine. Each entry states its
scope in one sentence; the material following documents coverage.

\subsection{Kinematics}

\textbf{Motion under constant acceleration, treated without reference to its causes.}

Five equations describe all such motion:
\begin{align}
v   &= v_0 + at \\
x   &= x_0 + v_0 t + \tfrac{1}{2}at^2 \\
v^2 &= v_0^2 + 2a(x - x_0) \\
x   &= x_0 + \tfrac{1}{2}(v + v_0)t \\
x   &= x_0 + vt - \tfrac{1}{2}at^2
\end{align}

Each omits exactly one of the six quantities $\{x_0, x, v_0, v, a, t\}$, so any three knowns
determine the rest. The solver selects the appropriate equation and reports the derivation steps ---
the module is written for instructional use.

Also implemented: one-dimensional motion (uniform, uniformly accelerated, free fall, vertical throw);
two-dimensional motion by independent component decomposition; projectile motion with range, apex,
flight time and impact velocity; and relative motion via $v_{AC} = v_{AB} + v_{BC}$.

\subsection{Rigid body dynamics}
\label{sec:rigidbody}

\textbf{Bodies that rotate without deforming --- gyroscopes, spinning tops, spacecraft.}

Orientation is represented in one of two ways:

\begin{itemize}[leftmargin=1.4em,itemsep=0.3em]
  \item \textbf{Euler angles} $(\phi, \theta, \psi)$: precession, nutation and spin, in the ZYZ
        convention. Geometrically interpretable, but subject to gimbal lock at $\theta = 0, \pi$,
        where two axes align and the parameterization degenerates.
  \item \textbf{Quaternions} $(q_0, q_1, q_2, q_3)$: four parameters subject to the unit-norm
        constraint $q_0^2 + q_1^2 + q_2^2 + q_3^2 = 1$. Singularity-free at every orientation, at the
        cost of one redundant parameter and less direct interpretation.
\end{itemize}

The symmetric top has Lagrangian
\begin{equation}
L = \tfrac{1}{2}I_1\!\left(\dot\theta^2 + \dot\phi^2\sin^2\theta\right)
  + \tfrac{1}{2}I_3\!\left(\dot\psi + \dot\phi\cos\theta\right)^2 - Mgl\cos\theta
\end{equation}

Rotational kinetic energy is $\tfrac{1}{2}\omega^{\mathsf{T}} I \omega$, where $I$ is the inertia
tensor --- the rotational analogue of mass, encoding the distribution of mass about the rotation
axes. Angular velocity from quaternion rates is $\omega = 2E\dot q$.

\subsection{Central forces and orbital mechanics}

\textbf{Motion under a force depending only on separation --- gravitation and electrostatics.}

For $V(r)$ the Lagrangian is $\tfrac{1}{2}m(\dot r^2 + r^2\dot\phi^2) - V(r)$. The angle $\phi$ is
cyclic, so angular momentum $L = mr^2\dot\phi$ is conserved (Section~\ref{sec:noether}), and the
two-dimensional problem reduces to one-dimensional radial motion in the \textbf{effective potential}
\begin{equation}
V_{\mathrm{eff}}(r) = V(r) + \frac{L^2}{2mr^2}
\end{equation}
The second term is the \textbf{centrifugal barrier}, representing the angular momentum that must be
overcome to approach the centre.

Turning points satisfy $E = V_{\mathrm{eff}}(r)$, and their number and arrangement classify the orbit
as circular, elliptical, parabolic or hyperbolic. Includes the Kepler problem and apsidal precession
rates for non-inverse-square potentials.

\subsection{Collisions and scattering}

\textbf{Impulsive interactions and deflection by a potential.}

Momentum is conserved in all collisions; kinetic energy only in elastic ones. The \textbf{coefficient
of restitution} $e$ parameterizes the range: $e = 1$ elastic, $0 < e < 1$ inelastic, $e = 0$
perfectly inelastic with the bodies coalescing. Includes centre-of-mass frame transformations and
impulse calculations.

Scattering treats a particle incident from infinity, deflected through angle $\theta$, departing to
infinity. Covers the impact-parameter-to-angle relation, differential and total cross-sections, and
Rutherford scattering.

\subsection{Variable-mass systems}

\textbf{Systems whose mass changes during motion --- principally rockets.}

Newton's second law in the form $F = ma$ does not apply. The correct statement retains the momentum
flux term:
\begin{equation}
F_{\mathrm{ext}} = m\frac{\ud v}{\ud t} + v_{\mathrm{rel}}\frac{\ud m}{\ud t}
\end{equation}
The second term accounts for momentum carried by mass entering or leaving the system. A rocket
accelerates by expelling mass, not by reaction against an external medium. Includes the Tsiolkovsky
equation, ejection and accretion, conveyor belts and chain dynamics.

\subsection{Continuum mechanics and classical fields}

\textbf{Systems with continuously distributed degrees of freedom --- strings, membranes, fields.}

The configuration is a field $\phi(x, t)$ rather than finitely many coordinates, and the action
becomes a double integral over space and time. The Euler--Lagrange equation generalizes to
\begin{equation}
\dd{\mathcal{L}}{\phi}
  - \frac{\partial}{\partial t}\dd{\mathcal{L}}{\phi_t}
  - \frac{\partial}{\partial x}\dd{\mathcal{L}}{\phi_x} = 0
\end{equation}
yielding the wave equation. Covers vibrating strings, membrane vibration and the stress--energy
tensor.

\subsection{Fluid dynamics}

\textbf{Free-surface and splashing flows, by particle methods.}

\textbf{Smoothed Particle Hydrodynamics} discretizes a fluid into particles whose properties are
distributed over a finite smoothing radius $h$ rather than concentrated at points. Field values at
any location are recovered by weighted summation over neighbours within that radius, which bounds the
cost per particle. The method handles free surfaces and topological changes naturally, where grid
methods require special treatment.

Three kernels are used:

\begin{center}
\begin{tabular}{@{}lcl@{}}
\toprule
\textbf{Kernel} & \textbf{Form} & \textbf{Application} \\
\midrule
Poly6                & $\dfrac{315}{64\pi h^9}\left(h^2 - r^2\right)^3$ & Density \\[0.9em]
Spiky gradient       & $-\dfrac{45}{\pi h^6}(h - r)^2\,\hat{r}$         & Pressure force \\[0.9em]
Viscosity Laplacian  & $\dfrac{45}{\pi h^6}(h - r)$                     & Viscous force \\
\bottomrule
\end{tabular}
\end{center}

Density is $\rho_i = \sum_j m_j W(r_{ij}, h)$, with pressure from a Tait-form equation of state
$p = k(\rho - \rho_0)$. Forces comprise pressure (pairwise symmetrized), viscosity and gravity;
integration is semi-implicit Euler. Boundary conditions include no-slip and free-slip walls, periodic
and open.

\subsection{Electromagnetism}

\textbf{Charged particle dynamics, wave propagation and plasma behaviour.}

The charged-particle Lagrangian incorporates the scalar and vector potentials,
\begin{equation}
L = \tfrac{1}{2}mv^2 - q\varphi + q\,\vec{v}\cdot\vec{A}
\end{equation}
reproducing the \textbf{Lorentz force} $\vec F = q(\vec E + \vec v \times \vec B)$.

The conjugate momentum is $\vec p = m\vec v + q\vec A$, not $m\vec v$ --- a direct illustration that
momentum in this formalism is defined by $\partial L/\partial\dot q$ and need not coincide with the
elementary definition.

Also implemented: cyclotron frequency and Larmor radius; magnetic mirrors with loss cones; Penning
traps (cyclotron, axial, magnetron and modified-cyclotron frequencies, and the invariance theorem);
guiding-centre drifts (grad-$B$, curvature, polarization, $\vec E \times \vec B$); plasma frequency
and Debye length; electromagnetic waves (impedance, intensity, energy density, radiation pressure);
antennas (Hertzian and half-wave dipoles, gain, effective aperture); rectangular waveguides (cutoff
frequencies, guide wavelength, phase and group velocity, TE/TM impedance); skin effect; and the
Biot--Savart law.

\subsection{Special relativity}

\textbf{Dynamics at speeds comparable to $c$.}

The relativistic Lagrangian $L = -mc^2\sqrt{1 - v^2/c^2} - V$ is equivalent to extremizing proper time
along the worldline: of all paths between two events, the physical one maximizes the elapsed time
measured by the traveller.

The \textbf{Lorentz factor} $\gamma = 1/\sqrt{1 - v^2/c^2}$ governs all relativistic corrections,
approaching $1$ at low speed and diverging as $v \to c$. Core relations: $p = \gamma mv$,
$E = \gamma mc^2$, $T = (\gamma - 1)mc^2$, and the invariant $E^2 = (pc)^2 + (mc^2)^2$.

Four-vectors are classified as timelike, spacelike or lightlike by the sign of their invariant. Also
covers Lorentz boosts, relativistic velocity addition, time dilation, length contraction, rapidity,
relativistic collisions and threshold energies, Mandelstam $s$, longitudinal and transverse Doppler
shift, aberration, Compton scattering, synchrotron radiation, Thomas precession, and the twin paradox
for both constant velocity and constant proper acceleration.

\subsection{General relativity}

\textbf{Gravitation as spacetime curvature; black holes and cosmology.}

\textbf{Schwarzschild geometry} (static, spherically symmetric):
\begin{equation}
\ud s^2 = -\left(1 - \frac{r_s}{r}\right)c^2\ud t^2
        + \left(1 - \frac{r_s}{r}\right)^{-1}\!\ud r^2 + r^2\,\ud\Omega^2
\end{equation}
with $r_s = 2GM/c^2$. Provides gravitational redshift, the photon sphere at $1.5\,r_s$, the innermost
stable circular orbit at $3\,r_s$, tidal acceleration, surface gravity and Hawking temperature.
Geodesics --- the trajectories of free particles, which are the straightest available paths in curved
geometry --- are integrated from an effective potential, with separate treatment for massive and
massless particles.

\textbf{Kerr geometry} (rotating): inner and outer horizons, the ergosphere (within which frame
dragging makes static observers impossible), frame-dragging rate, and spin-dependent prograde and
retrograde ISCO radii.

\textbf{Gravitational lensing}: deflection $\alpha = 4GM/(c^2 b)$, Einstein radius, magnification and
time delay.

\textbf{FLRW cosmology}: Hubble parameter as a function of redshift, Hubble time and distance,
comoving, luminosity and angular-diameter distances, lookback time, cosmic age, critical density and
the deceleration parameter.

\subsection{Quantum mechanics}

\textbf{Semiclassical methods --- the regime where classical quantities determine quantum results.}

The \textbf{WKB approximation} treats the wavefunction as a slowly modulated wave with local
wavenumber set by the classical momentum $p(x) = \sqrt{2m(E - V(x))}$. The \textbf{Bohr--Sommerfeld
condition} then quantizes the classical action:
\begin{equation}
\oint p\,\ud x = \left(n + \tfrac{1}{2}\right)h
\end{equation}
Only orbits whose enclosed action is a half-integer multiple of Planck's constant are permitted,
which is the origin of discrete energy levels. This is the point at which the classical machinery of
the rest of the library yields quantum results.

Exactly solved systems: the harmonic oscillator ($E_n = \hbar\omega(n + \tfrac{1}{2})$, Hermite
wavefunctions, uncertainty products); infinite and finite square wells; step potentials;
delta-function barriers; and the hydrogen atom ($E_n = -13.6\ \mathrm{eV}/n^2$, Bohr radii,
degeneracy, transition wavelengths, spectral series).

\textbf{Tunneling} --- transmission through a barrier exceeding the particle's energy, classically
forbidden --- covers rectangular barriers, WKB transmission, the Gamow factor, alpha decay rates,
double-well splitting and resonant tunneling. \textbf{Ehrenfest's theorem} propagates expectation
values classically, establishing the correspondence limit.

\subsection{Thermodynamics}

\textbf{Heat engines, equations of state and phase behaviour.}

\textbf{Carnot efficiency} $\eta = 1 - T_c/T_h$ bounds every heat engine operating between two
reservoirs --- a consequence of the second law rather than an engineering limitation, and the reason
thermal power generation necessarily rejects heat.

Also implemented: the Otto cycle ($\eta = 1 - r^{1-\gamma}$, monotonic in compression ratio), Diesel
and Rankine cycles; the van der Waals equation
$\left(P + an^2/V^2\right)(V - nb) = nRT$ with critical point and compressibility factor;
Redlich--Kwong; phase transitions via Clausius--Clapeyron; Debye and Einstein heat capacity models;
and the four Maxwell relations.

\subsection{Statistical mechanics}

\textbf{Derivation of macroscopic thermodynamics from microscopic states.}

The \textbf{Boltzmann distribution} assigns probability proportional to $e^{-E/k_B T}$, with the
\textbf{partition function} $Z = \sum_i e^{-E_i/k_B T}$ as the generating quantity from which
thermodynamic properties follow.

Includes the Maxwell speed distribution with most-probable, mean and RMS speeds; ideal gas pressure,
internal energy, heat capacities, Sackur--Tetrode entropy and free energies; quantum oscillator
ensembles; the Ising model with Metropolis Monte Carlo and the exact two-dimensional critical
temperature; and Fermi--Dirac and Bose--Einstein occupation statistics with Fermi energy and BEC
critical temperature.

\subsection{Solid mechanics}

\textbf{Deformation and failure of engineering materials.}

\begin{itemize}[leftmargin=1.4em,itemsep=0.35em]
  \item \textbf{Elasticity} --- Hooke's law for isotropic, transversely isotropic, orthotropic,
        monoclinic and triclinic symmetry classes.
  \item \textbf{Stress and strain} --- principal stresses, invariants $I_1, I_2, I_3, J_2, J_3$, von
        Mises and Tresca yield criteria, Mohr's circles in two and three dimensions, tensor
        transformations, stress concentration factors and strain rosette reduction.
  \item \textbf{Beams} --- Euler--Bernoulli theory for slender beams and Timoshenko theory where
        shear deformation is significant; deflection, slope, bending moment and shear; composite,
        tapered and curved beams.
  \item \textbf{Plates} --- Kirchhoff--Love for thin plates and Mindlin--Reissner for thick;
        rectangular and circular solutions, buckling and natural frequencies.
  \item \textbf{Fracture} --- stress intensity factors for Modes I, II and III; energy release rate;
        the $J$-integral; CTOD; the Griffith criterion (propagation when released strain energy
        exceeds the surface energy cost); Paris law for fatigue crack growth; and plastic zone size.
  \item \textbf{Fatigue} --- failure under cyclic loading below the static strength. S--N curves
        $\sigma_a = \sigma_f'(2N_f)^b$, strain-life, mean stress correction and cumulative damage.
  \item \textbf{Viscoelasticity} --- Maxwell, Kelvin--Voigt and Standard Linear Solid models;
        generalized Maxwell via Prony series; creep and relaxation functions; complex modulus;
        time--temperature superposition; and Boltzmann superposition.
  \item \textbf{Thermal stress} --- $\varepsilon = \alpha\Delta T$ linear and $3\alpha\Delta T$
        volumetric, with thermoelastic coupling. Constrained thermal expansion generates stress
        independent of applied load.
\end{itemize}

\subsection{Optics}

\textbf{Propagation, imaging and interference of light.}

Geometric optics (ray tracing, Snell's law, reflection, thin and thick lenses), wave optics
(interference and diffraction), polarization via Jones matrices and Stokes parameters, and fibre
optics.

\subsection{Acoustics}

\textbf{Generation, propagation and confinement of sound.}

Wave equation solutions in one, two and three dimensions; resonance and standing waves; acoustic
impedance and boundary reflection; the Doppler effect; room acoustics and Sabine reverberation; and
ultrasonics.

% ============================================================================
\section{Summary}
% ============================================================================

The library rests on a single organizing principle:

\begin{quote}
\itshape
A mechanical system is specified by one scalar function --- kinetic energy minus potential energy ---
and everything else follows by differentiation.
\end{quote}

Given $L = T - V$, a fixed procedure yields the equations of motion; inverting the mass matrix yields
accelerations; a Legendre transform yields the energy; absent coordinates yield conservation laws;
structure-preserving integrators yield trajectories that respect those conservation laws over
arbitrarily long runs; and the same symbolic expressions compile to thirteen target languages.

The Newtonian formulation yields identical predictions. It requires, however, that a person derive
them --- once per system, once per coordinate choice, by hand, correctly.

That is precisely the step a compiler cannot perform on the user's behalf, and it is why the library
is built on the other formulation.

\end{document}
